% Source: https://artofproblemsolving.com/community/c3h2031652_algebra
% Stable identity: aops-topic-2031652-q01
\textbf{Case $a=0$.} The equation is $3x+3=0$, with the single root $x=-1$, an integer. So $a=0$ works.

\textbf{Case $a\ne0$.} The equation is quadratic. Suppose its roots $r,s$ (counted with multiplicity) are integers. By Vieta's formulas,
\[
r+s=-\dfrac{3+4a}{a}=-4-\dfrac{3}{a},\qquad rs=\dfrac{2a^2+4a+3}{a}=2a+4+\dfrac{3}{a}.
\]
The first shows $m=\dfrac{3}{a}$ is a (nonzero) integer; then the second shows $2a=\dfrac{6}{m}$ is an integer, so $m$ divides $6$. Thus
\[a\in\left\{\pm3,\ \pm\dfrac{3}{2},\ \pm1,\ \pm\dfrac{1}{2}\right\}.\]
The discriminant is
\[
\Delta=(3+4a)^2-4a(2a^2+4a+3)=9+12a-8a^3.
\]
Checking the eight candidates:
\[
\begin{array}{c|cccccccc}
a & 3 & -3 & \dfrac{3}{2} & -\dfrac{3}{2} & 1 & -1 & \dfrac{1}{2} & -\dfrac{1}{2}\\[4pt]\hline
\Delta & -171 & 189 & 0 & 18 & 13 & 5 & 14 & 4
\end{array}
\]
Integer roots need $\Delta$ to be a perfect square, leaving $a=\dfrac{3}{2}$ and $a=-\dfrac{1}{2}$.
\begin{itemize}
\item $a=\dfrac{3}{2}$: the equation is $\dfrac{3}{2}x^2+9x+\dfrac{27}{2}=\dfrac{3}{2}(x+3)^2=0$, with the double root $x=-3$.
\item $a=-\dfrac{1}{2}$: the equation is $-\dfrac{1}{2}x^2+x+\dfrac{3}{2}=-\dfrac{1}{2}(x-3)(x+1)=0$, with roots $3$ and $-1$.
\end{itemize}
Both work. Therefore
\[\boxed{a\in\left\{-\dfrac{1}{2},\ 0,\ \dfrac{3}{2}\right\}}.\]
